% !TEX root = rosary.tex
% !TEX recipe = LuaLaTeX+se
%\documentclass{article}
\documentclass[parskip]{scrartcl} % set document class: manual at https://ctan.org/pkg/koma-script
%\usepackage{fontspec}

% Load packages:
\usepackage[osf,p]{libertine} % set font
\usepackage[autocompile]{gregoriotex} % enable Gregorio score inclusion
\usepackage[latin]{babel} % set language
\usepackage{multicol}
\usepackage[margin=3cm]{geometry}
\usepackage{scrlayer-scrpage}
\usepackage{musixtex}

\newcommand{\myhome}{C:/Users/johna/Chant/}
\gresetgregpath{{\myhome}}

\setkomafont{section}{\normalfont\centering\huge\scshape} % section heading style
\setkomafont{subsection}{\normalfont\centering\Large\scshape} % section heading style
\setkomafont{subsubsection}{\normalfont\centering\small\scshape} % section heading style
\setcounter{secnumdepth}{-\maxdimen} % remove section numbering

\setkomafont{pagehead}{\scshape}
\automark[section]{section}
%\automark*[subsection]{subsection}

%\lohead[\thesection]
%{\thesection}
%\rohead[\thesubsection]
%{\thesubsection}
\pagestyle{scrheadings}

\def\translation{NABRE}

% Set the space around the initial:
% See http://gregorio-project.github.io/gregoriotex/details.html for more details and options
\grechangedim{beforeinitialshift}{2.2mm}{scalable}
\grechangedim{afterinitialshift}{2.2mm}{scalable}

% Set the initial font (change 43 for a larger size):
%\grechangestyle{initial}{\fontsize{43}{43}\selectfont}%

% Make staff lines red; remove for black:
%\gresetlinecolor{gregoriocolor}

% Use the "commentary" field of the score in the top right corner:
\gresetheadercapture{commentary}{grecommentary}{string}

% Format annotation above initial
\grechangestyle{annotation}{\small\bfseries}

\begin{document}

\thispagestyle{plain}
\section*{Foreword}

\subsection{``For feasts?'' ``Ferias?'' What?}
The Church observes three ranks of feast days throughout the year:
\textbf{memorials}, \textbf{feasts}, and \textbf{solemnities}.
Solemnities are usually Holy Days of Obligation, but also include
your own parish's patronal feast and some local
observances. A \textbf{feria} is a normal weekday on which no saint or
other feast is commemorated.

We often forget that, even as we leave Sunday Mass and begin our week,
the Church continues to pray and celebrate God's graces without ceasing,
as the Apostle Paul commands us. 

Multiple options are given to sing, both for the mysteries and for the prayers and orations
surrounding the mysteries. Often, these are labeled with \emph{suggested}
use cases. This author has no authority whatsoever to do anything more 
than suggest anything---so please feel free to use your own discretion.

Beginners, in fact, should only sing the simpler tones
(regardless of liturgical season) until comfortable.
However, experienced singers should likewise not ignore the simpler tones!
Reserving the more ornate solemn tones for special days is a wonderful
way of uniting oneself to the life of the Church.



\subsection{Reading Chant}
If you are brand-new to chant, please don't get too lost in the weeds 
here---you might even skip this section altogether and be fine. 
Even if you are singing with others, as long as everyone
sings at the same cadence, you will be singing mostly the same
notes and it will still be a prayerful experience.

\subsubsection{Relative pitch}
Even if you don't know what it means, you probably have seen 
the ``G'' looking figure at the beginning of this staff:
\begin{music}
    \startextract
    \Notes\ibbbu0h0\qb0e\tbbbu0\qb0e\tbbu0\qb0e\tbu0\qb0e\enotes
    \zendextract
\end{music}

\subsubsection{The staff/clef}
Unlike modern notation, the staff for Gregorian chant has only 
four lines instead of five---this reflects the much smaller range of
the voice compared to most other instruments.

The \textbf{clef} is a symbol at the beginning of a staff that indicates
the relative position of notes. 

Note that the clef kind of ``circles'' the second line from the bottom.
This particular clef communicates an absolute pitch---on most instruments,
392Hz, or G.


\subsubsection{Reciting tones}
A \textbf{reciting tone} is less of a chant in and of itself,
and more of a ``formula'' for chanting arbitrary phrases.
Prayers sung according to a reciting tone will be sung \emph{mostly}
on the same note, but deviate at the beginning, end, and around commas 
or pauses.

The way these typically work is by using a \textbf{hollow note} 
to denote places where a pitch can be held for an arbitrary
number of syllables, as needed to match the end of the tone to the
end of the sentence or phrase.

\textbf{Accent marks} in the reciting tone are used as a cue for
how to match the ending part of the tone to the end of the sentence.
They should typically line up with the stressed syllable in the word 
being sung.

For example,
\gabcsnippet{(c3)The(h) ...(hr) ...(hr) my(h)ste(f)ry:(f) (;) ...(h hr) (g) (f) ...(hrr1 h.) (::)}
becomes
\gabcsnippet{(c3)The(h) first(h) sor(h)row(h)ful(h) my(h)ste(f)ry:(f) (;) 
the(h) a(h)gon(h)y(h) in(g) the(f) gar(h)den.(h.) (::)}
The accent over the hollow \textbf{do}
indicates that the last stressed syllable of the sentence is to be 
sung on \textbf{do}---when this is the last (or only) syllable of the word,
only one \textbf{do} is sung after the \textbf{ti-la} after the mediant.
So the fourth sorrowful mystery is sung:
\gabcsnippet{(c3)The(h) fourth(h) sor(h)row(h)ful(h) my(h)ste(f)ry:(f) (;) 
the(h) car(h)ry(h)ing(h) of(g) the(f) Cross.(h.) (::)}
``Cross'' has only one syllable, so it by default is the stressed syllable,
and so the tone returns to \textbf{do} on ``Cross''.
Since ``garden'' is stressed in its first syllable, the tone returns to 
\textbf{do} on ``\textbf{gar}''.


\section{The texts of the prayers of the Rosary}
\subsection{Overall structure of the Rosary}
\paragraph{The Mysteries.}
There are four sets of mysteries: Joyful, Luminous, Sorrowful,
and Glorious. The Luminous Mysteries are an optional recent addition
to the Rosary. They began as a separate devotion, and were more
recently added to the Rosary by Pope St. John Paul II.

If praying the Luminous, it is recommended to pray the Joyful on
Monday and Saturday; the Sorrowful Tuesday and Friday; the Luminous Thursday;
the Glorious Wednesday and Sunday.

If one does not pray the Luminous mysteries, the traditional rotation
is Joyful on Monday and Thursday, Sorrowful on Tuesday and Friday, 
Glorious on Wednesday and Saturday. (A benefit of this rotation is
that the mysteries are prayed in order.)

On Sunday, the Joyful Mysteries
are prayed during the Advent and Christmas season; the Sorrowful during 
Lent; and Glorious otherwise.

The Mysteries are:
\begin{multicols}{2}
    \begin{itemize}
        \item \textbf{Joyful}
        \item[1] The Annunciation of Gabriel to Mary
        \item[2] The Visitation of Mary to Elizabeth
        \item[3] The Birth of Jesus
        \item[4] The Presentation in the Temple
        \item[5] The Finding of the Child Jesus in the Temple
    \end{itemize}
    \columnbreak
    \begin{itemize}
        \item \textbf{Luminous}
        \item[1] The Baptism of the Lord
        \item[2] The Wedding Feast at Cana
        \item[3] The Proclamation of the Kingdom
        \item[4] The Transfiguration
        \item[5] The Institution of the Holy Eucharist
    \end{itemize}
\end{multicols}
\begin{multicols}{2}
    \begin{itemize}
        \item \textbf{Sorrowful}
        \item[1] The Agony in the Garden
        \item[2] The Scourging at the Pillar
        \item[3] The Crowning with Thorns
        \item[4] The Carrying of the Cross
        \item[5] The Crucifixion and Death of Jesus
    \end{itemize}
    \columnbreak
    \begin{itemize}
        \item \textbf{Glorious}
        \item[1] The Resurrection of Jesus from the Dead
        \item[2] The Ascension of Jesus into Heaven
        \item[3] The Descent of the Holy Spirit
        \item[4] The Assumption of Mary
        \item[5] The Coronation of Mary as Queen of Heaven and Earth
    \end{itemize}
\end{multicols}


\paragraph{Introduction.} 
After deciding which set of mysteries are to be prayed---or all four
(or three) of them!---one begins:
\begin{itemize}
    \item Begin with the Apostle's Creed. \emph{on the Crucifix}
    \item One Our Father. \emph{On the large bead next to the Crucifix}
    \item Three Hail Marys \emph{On the three small beads following the first Our Father bead.}
    \item One Glory Be. \emph{On the chain separating the last Hail Mary bead from the next Our Father bead.}
\end{itemize}

\paragraph{The First Mystery.}
\begin{itemize}
    \item Announce the first Mystery.
    \item One Our Father. \emph{On the Our Father bead}
    \item A short verse can be recited to call to mind a
        particular aspect of this Mystery---or, as we will learn
        in this booklet, an antiphon pertaining to this Mystery is sung.
    \item Ten Hail Marys. \emph{On the ten smaller beads}
    \item One Glory Be. \emph{On the chain separating the last Hail Mary from the next decade.}
    \item If singing in the manner proposed in this booklet,
        the antiphon sung before the first Hail Mary is repeated.
\end{itemize}

\paragraph{Subsequent Mysteries.}
The steps above are repeated until all Mysteries in the set are
completed.

\paragraph{Conclusion.}


\section{Recto tono}
\emph{Recto tono} means ``straight tone'', and this 
simply means to sing everything on the same pitch.
This is the easiest way to chant any prayer. Simply pick a pitch 
that is comfortable for you, and sing, somewhat slowly, and 
all on the same note:
\gabcsnippet{(c3)The(h) first(h) sor(h)row(h)ful(h) my(h)ste(h)ry:(h) (;) 
the(h) a(h)gon(h)y(h) in(h) the(h) gar(h)den.(h.) (::)}

Pause for a brief moment, and then continue with the Our Father:
\gabcsnippet{(c3)Our(h) Fa(h)ther,(h) who(h) art(h) in(h) hea(h)ven,(h)
(,) ...(hr) (::)}

This is a great way for beginner singers to get started, but 
even experienced singers should make use of this tone for penitential
days, such as Fridays and weekdays in Lent---or, even just Fridays in Lent.

One can proceed with the mysteries also \emph{recto tono}, or use the 
mysteries given in the following section, Simple Tones.

After each mystery, sing the \emph{O my Jesus} prayer \emph{recto tono}:
\gabcsnippet{(c3)O(h) my(h) Je(h)sus(h), for(h)give(h) us(h) our(h) sins...(h)}


%% everything before here is too wordy
\clearpage
\section{Simple Tones}
For ferias, memorials, and Sundays in Lent.
\gabcsnippet{(c3)The(h) ...(hr) ...(hr) my(h)ste(f)ry:(f) (;) 
the(h) a(h)gon(h)y(h) in(g) the(f) gar(hr1)den.(h.) (::)}
or,
\gresetinitiallines{0}
\gabcsnippet{(c3) 
...the(h) birth(g) of(f) Je(hr1)sus.(h.) (::)}
or,
\gabcsnippet{(c3) 
...the(h) crow(h)ning(g) of(f) thorns.(h.) (::)}
or,
\gabcsnippet{(c3) 
...the(h) Res(h)sur(h)rec(h)tion(h) of(h) Je(h)sus(h) from(g) the(f) dead.(h.) (::)}
\gresetinitiallines{1}

After the mystery is announced, pause for a brief moment, then sing
the Our Father:
\gregorioscore{Communes/pater-noster/English-ferialis}


%\begin{center}
%\greseparator{2}{30}
%\end{center}


\subsection{Joyful Mysteries}

\subsubsection{First Mystery - The Annunciation of Gabriel to Mary}
\gregorioscore{rosary/Joyful/I/theangelof}

\subsubsection{Fourth Mystery - The Presentation in the Temple}
\gregorioscore{rosary/Joyful/IV/viderunt-oculi}


\subsection{Sorrowful Mysteries}


\subsubsection{First Mystery - The Agony in the Garden}
\gregorioscore{rosary/Sorrowful/I/deusmeuseripeme_english}

\subsubsection{Second Mystery - The Scourging}
\gregorioscore{rosary/Sorrowful/II/trulyhebore}

\subsubsection{Third Mystery - The Crowning of Thorns}
\gregorioscore{rosary/Sorrowful/III/mykingdom}

\subsubsection{Fourth Mystery - The Carrying of the Cross}
\gregorioscore{rosary/Sorrowful/IV/quivult}

%\gresetinitiallines{0} 
%\gregorioscore{HailMary8c}
%\gresetinitiallines{1} 

\subsection{Glorious Mysteries}

\subsubsection{Second Mystery - The Ascension of the Lord}
\gregorioscore{rosary/Glorious/II/goforth}


\subsubsection{Fourth Mystery - the Assumption of Mary}
%\gregorioscore{rosary/agreatsign}

\subsubsection{Fifth Mystery - the Coronation}
\gregorioscore{rosary/Glorious/V/thekingrose}

%\gresetinitiallines{0} 
\gregorioscore{rosary/Aves/HailMary8c}
%\gresetinitiallines{1} 

\clearpage
\section{Solemn tones}

For feasts, solemnities, and Sundays outside of Lent; also may 
be used the evening \emph{before} a solemnity or Sunday outside of Lent.

Announce the mystery using this formula:
\gabcsnippet{(c4)The(g) first (h)<i>second, third...</i>(hr) sor(h)row(h)ful(h hr) my(hr1)ste(g)ry:(g) (;) 
the(g) a(h)go(h)ny(h) in(g) the(g) gar(hr1)den.(h.) (::)}
or,
\gresetinitiallines{0}
\gabcsnippet{(c3) 
...the(h) birth(g) of(g) Je(hr1)sus.(h.) (::)}
or,
\gabcsnippet{(c3) 
...the(h) crow(h)ning(g) of(g) thorns.(h.) (::)}
or,
\gabcsnippet{(c3) 
...the(g) Res(h)sur(h)rec(h)tion(h) of(h) Je(h)sus(h) from(g) the(g) dead.(h.) (::)}
\gresetinitiallines{1}

After the mystery is announced, pause for a brief moment, then sing
the Our Father:
\gregorioscore{Communes/pater-noster/English-C}\label{pater:solemnis}
\begin{center}
\greseparator{2}{30}
\end{center}


\subsection{Joyful Mysteries}

\subsubsection{First Mystery - the Annunciation}
\gabcsnippet{(c4)The(g) first (h) joy(h)ful(h) my(h)ste(g)ry:(g) (;) 
the(g) an(h)nun(h)ci(h)a(h)tion(h) of(h) Ga(h)bri(h)el(g) to(g) Ma(h)ry.(h.) (::)}

\emph{Our Father, pg.~\ref{pater:solemnis}}

\gregorioscore{rosary/Joyful/I/ne-timeas-maria}
\gregorioscore{rosary/Aves/HailMary8c}

\begin{center}
    \greseparator{3}{20}
\end{center}

\subsection{Sorrowful Mysteries}

\subsubsection{Second Mystery - The Scourging}
\gregorioscore{rosary/Sorrowful/II/vere-languores}

\subsection{Gloroius Mysteries}

\subsubsection{Fourth Mystery - the Assumption of Mary}
%\gregorioscore{rosary/agreatsign}


\end{document}